\par\noindent\begin{minipage}{\linewidth}
% !TeX root = ../教材题库_全部条目_转写工作稿.tex
% Source: XJ-B1-C04-HW-P113-Y1566
% 第4章 / 4.2 指数函数 / 教材 p.113、114 / PDF 第120、121页
% 已对照原页转写题干及公式；待教师复核；未制作答案
% 原题及相邻原解截图（核对用）：../题目原图/XJ-B1-C04-HW-P113-Y1566-01.webp, ../题目原图/XJ-B1-C04-HW-P113-Y1566-02.webp
% 本文件只在转写完成后作为题干使用；OCR 条目尚不是可靠题干。
\begin{bankitem}{习题4.2 第6题}{第4章 / 4.2 指数函数 / 教材 p.113、114 / PDF 第120、121页}{XJ-B1-C04-HW-P113-Y1566}{已对照原页转写题干及公式；待教师复核；未制作答案}
在同一平面直角坐标系中画出 $y=4^x$，$y=4^{-x}$，$y=4^{x+1}$，$y=4^{x-1}$ 的图象，并说明后三个函数的图象分别可以由第一个函数的图象怎样变换得到。
\end{bankitem}
\end{minipage}\par\medskip
